
\textbf{Corollaire 4.15.} \uline{Soit} @@M0@@ \uline{le} @@M1@@-\uline{anneau engendré sur} @@M2@@ \uline{par l'indéterminée} @@M3@@ \uline{soumise aux relations} @@M4@@ \uline{pour} @@M5@@, \uline{i}.\uline{e}. @@M6@@ (4.9). \uline{Alors la filtration de} @@M7@@-\uline{anneau de} @@M8@@ \uline{est identique à la filtration par le poids par rapport aux} @@M9@@.

Compte tenu de la proposition 4.12, de la remarque 4.13.ii) et du fait que @@M10@@, @@M11@@, donc @@M12@@ est l'idéal engendré par les monômes de poids @@M13@@ par rapport aux @@M14@@.

\section*{5. Les @@M15@@.}

Ce paragraphe ne sera utilisé que dans Exp.VII~4, pour l'étude des immersions régulières~; il peut donc être omis en première lecture.

5.1. Soit @@M16@@ un entier @@M17@@. On désigne par @@M18@@ la @@M19@@-@@M20@@-algèbre engendrée par les générateurs @@M21@@, @@M22@@ soumis aux relations @@M23@@ pour @@M24@@~; on a donc :
@@M25@@

On pose @@M26@@.

\textbf{Lemme 5.2.} \uline{Pour tout} @@M27@@, @@M28@@ \uline{est divisible par} @@M29@@.

On sait (2.3) que @@M30@@ est un polynôme isobare de poids @@M31@@ par rapport aux @@M32@@~; il suffit donc de montrer le lemme pour les @@M33@@.

Soit @@M34@@ le @@M35@@-anneau engendré sur @@M36@@ par les indéterminées @@M37@@, soumises aux relations @@M38@@ si @@M39@@. Il est clair que @@M40@@ pour @@M41@@. On peut donc définir un @@M42@@-homomorphisme @@M43@@ de @@M44@@ dans @@M45@@, en envoyant @@M46@@ sur @@M47@@~; on a alors @@M48@@,
